\section{Evaluation}
\label{sec:eval}

\subsection{Experimental Setup}

\textbf{Hardware.}  All measurements are on an NVIDIA RTX~4080~SUPER
(16\,GB VRAM) with an AMD Ryzen~9 7950X CPU.

\textbf{Pipeline.}  WorldFusion (Where2Comm attention fusion),
AB3DMOT tracking (vectorized Kalman), MTR prediction (25-step
multi-modal).  Post-backbone features are real, extracted from the
OPV2V~\cite{xu2022opencood} validation split through the trained
WorldFusion encoder (post-backbone, 256-channel, shrink-convolved).

\textbf{Agent scaling.}  $N \in \{4, 8, 16, 24, 32, 48, 64\}$.
For $N$ exceeding the number of unique OPV2V CAVs, features are
replicated with distinct transformations to preserve workload fidelity.

\textbf{Ablation configurations.}
Five policies isolate each mechanism:
(i)~\textit{static}: full pipeline, no adaptation;
(ii)~\textit{filter\_only}: input contribution filter with $K_{\max}{=}16$,
  no prediction adaptation;
(iii)~\textit{amort\_only}: divergence gating only, no filter, no risk budget;
(iv)~\textit{risk\_only}: risk-budgeted scheduling only (30\,ms prediction
  budget), no filter, no amortization;
(v)~\textit{adaptive\_all}: all three mechanisms enabled
  ($K_{\max}{=}16$, amortization, 30\,ms prediction budget).

\textbf{Ticks.}  30 ticks per configuration, with steady-state metrics
computed over the last 8 ticks (after tracker warmup).

\subsection{Per-Stage Latency Scaling}

\Cref{tab:stage_latency} reports per-stage latency for the static
pipeline.  Fusion scales linearly with $N$ (Where2Comm pairwise
attention), tracking scales linearly with the number of detected
objects $M$, and MTR prediction scales with the number of tracked
objects requiring inference.  Detection is near-constant ($\sim$1\,ms)
because it operates on the fused feature map.

The pipeline first exceeds 100\,ms at $N{=}16$ (149\,ms) and reaches
302\,ms at $N{=}32$.  Beyond $N{=}32$, timing is estimated
from the per-stage cost model (\cref{eq:cost}) calibrated on the
measured range.

\subsection{Five-Way Ablation}

\Cref{tab:ablation_timing,tab:ablation_deadline} present end-to-end
latency and deadline compliance for each configuration.

\textbf{Static baseline.}  Fails at $N{\geq}16$; latency grows to
184\,ms at $N{=}64$.  The non-monotonicity at $N{=}32$ (104\,ms)
reflects variation in detection count across ticks.

\textbf{Filter only.}  Capping agents at $K{=}16$ restores compliance
for $N{\geq}24$ but does not help at $N{=}16$ (where $K{=}N$).
The filter itself adds $<$2\,ms even at $N{=}64$.

\textbf{Amortization only.}  Reduces prediction cost by reusing cached
trajectories for non-divergent tracks.  Effective through $N{=}32$
(100\% compliance) but loses compliance at $N{=}48$ because fusion cost
alone approaches the deadline.

\textbf{Risk-budget only.}  Caps prediction at 30\,ms by serving only
the highest-risk divergent tracks.  Effective through $N{=}32$ but
fusion cost exceeds the budget at $N{=}48$.

\textbf{Adaptive (all three).}  The only configuration achieving
100\% deadline compliance at every $N$.  At $N{=}64$, end-to-end
latency is 50\,ms (50\% headroom), because the filter caps fusion
cost while amortization and risk budgeting cap prediction cost.

\subsection{Quality Metrics: Fixed-GT Evaluation}

To avoid the pitfalls of nearest-neighbor GT matching
(ghost predictions matched to arbitrary objects, missing predictions
silently ignored), we use a fixed-GT evaluation:
\begin{itemize}
  \item For each GT object, find the best matching prediction within
    a 2\,m radius.
  \item Unmatched GT objects contribute a 10\,m miss penalty to ADE.
  \item The miss rate (fraction of GT objects without a matching
    prediction) is reported explicitly.
  \item Quality is stratified by risk level: objects within 30\,m of
    the ego origin are \emph{high-risk}; others are \emph{low-risk}.
\end{itemize}

\begin{table}[t]
\caption{Quality metrics (steady-state, ticks $\geq 22$).  ADE includes
  a 10\,m penalty for missed GT objects.  High-risk objects are within
  30\,m of the ego origin.}
\label{tab:quality}
\small
\begin{tabularx}{\linewidth}{lCCCCC}
\toprule
\textbf{Config} & \textbf{N} & \textbf{ADE-all} & \textbf{ADE-hi}
  & \textbf{miss\%} & \textbf{miss-hi\%} \\
\midrule
\multicolumn{6}{c}{\textit{\todo{Numbers pending Multi-V2X evaluation}}} \\
\bottomrule
\end{tabularx}
\end{table}

\subsection{Predictor Selection: MTR over SMART}

\Cref{tab:predictor} compares MTR and SMART on the same OPV2V
evaluation set.  SMART's 137\,ms inference time alone exceeds the
100\,ms deadline.  Its 42.7\,m ADE (vs.\ MTR's 4.5\,m) reflects
a Waymo-to-CARLA domain gap.  MTR's single-call edge-mode inference
(56.8\,ms, 4.13\,m ADE) is statistically equivalent to the multi-call
CMP-style invocation (73.8\,ms, 4.10\,m ADE) at 1.3$\times$ lower
latency.

\begin{table}[t]
\caption{Predictor comparison on OPV2V data.}
\label{tab:predictor}
\small
\begin{tabularx}{\linewidth}{lCCC}
\toprule
\textbf{Model} & \textbf{ADE (m)} & \textbf{FDE (m)}
  & \textbf{Latency (ms)} \\
\midrule
SMART (Waymo-trained) & 42.70 & 50.25 & 137 \\
MTR (multi-call, CMP) &  4.10 &  9.00 &  74 \\
MTR (single-call, edge) & 4.13 & 8.85 &  57 \\
\bottomrule
\end{tabularx}
\end{table}

\subsection{V2V Networking Validation}

We run ns-3 5G-LENA NR~V2X Mode~2 sidelink simulations to validate
the analytical argument in \cref{sec:problem}.  At 200\,KB
intermediate fusion payloads, PRR${}=0$ at all $N$: the payload
exceeds the sidelink transport block size.  At 17\,KB compressed
payloads, PRR degrades from 0.30 ($N{=}2$) to 0.006 ($N{=}64$).
These results confirm that V2V intermediate fusion is unviable at
dense scale, independent of the MAC protocol or bandwidth allocation.

\subsection{Cost Model Validation}

\todo{Insert predicted-vs-measured latency scatter for each stage
  ($c_{\mathrm{fuse}}$, $c_{\mathrm{track}}$, $c_{\mathrm{pred}}$).
  Target: $R^2 > 0.95$ on held-out ticks.}

\subsection{Greedy Gap}

\todo{For $N \leq 8$ where ILP is tractable, report gap between
  greedy utility and oracle optimal.  Preliminary: gap $< 5$\% in
  all tested instances.}
